\newpage
\section*{Konklusjon}
\addcontentsline{toc}{section}{Konklusjon}

Forsøket undersøkte virkningsgraden til en Pelton-turbin ved to ulike nålposisjoner. Den høyeste målte virkningsgraden ble funnet ved nålposisjon 3, med 58,4\% ved et turbinturtall på 14,78 $s^{-1}$, tilsvarende omtrent 887 rpm. For nålposisjon 7 var maksimal virkningsgrad 50,7\% ved 14,67 $s^{-1}$, tilsvarende omtrent 880 rpm.

Resultatene viser at nålposisjonen påvirker hvor effektivt turbinen omdanner hydraulisk effekt til mekanisk effekt. Selv om nålposisjon 7 ga høyere volumstrøm og større hydraulisk effekt, ga den lavere maksimal virkningsgrad enn nålposisjon 3. Dette tyder på at nålposisjon 3 ga et mer gunstig driftspunkt i dette forsøket.

Siden bare to nålposisjoner ble undersøkt, kan forsøket ikke fastslå den optimale innstillingen for turbinen. For å gjøre dette mer presist bør flere nålposisjoner testes, og målingene bør gjentas for å redusere usikkerheten.
